\section{Angular soft data and fixed-order logarithmic recovery}
\label{sec:angularlogs}
Soft coefficients carry angular, helicity, and sometimes colour data.
They can be distributions rather than pointwise functions. A recovery
estimate for a prescribed angular pairing should therefore specify a
normed space and its dual. We give such a statement, including fixed-order
logarithmic terms, directly for Mellin observations.

Let $X$ be a nonzero complex Banach space. Angular examples are
$X=L^2(K,d\Omega)$ for a patch $K$ on the celestial sphere, or
$X=H^{-s}(S^2)$ paired with $H^s$ test functions. The field component
and hard kinematics are fixed. Membership in these spaces is an
assumption: collinear singularities need not belong to $L^2(S^2)$.
For fields with spin, the pairing uses an appropriate dual test section.
The angular norm here is distinct from the frequency-variable Sobolev
norm in Section~\ref{sec:boundary}.

For fixed $a,\alpha,M>0$ and an integer $m\ge0$, suppose
\begin{equation}\label{eq:vectorclass}
 \begin{gathered}
 f(x)=\sum_{k=0}^m u_k(\log(1/x))^k+R(x),\qquad
 \|R(x)\|_X\le Mx^\alpha\quad(x>0),\\
 Q_a(f)=\int_0^\infty x^{a-1}\|f(x)\|_X\,dx<\infty.
 \end{gathered}
\end{equation}
All integrals with values in $X$ are Bochner integrals, and the
functions are assumed strongly measurable. The remainder assumption
is global; it is not inferred from a finite observation window.
As before, $x$ is numerical energy in a fixed reference convention,
and a known leading power has been stripped off.
Splitting the Mellin integral at one shows that
\[
 G(z)=\int_0^\infty x^{z-1}f(x)\,dx
 \quad\hbox{has principal part}\quad
 \sum_{k=0}^m\frac{k!u_k}{z^{k+1}}.
\]
Indeed, the infrared remainder is holomorphic for $\Re z>-\alpha$,
and the ultraviolet integral for $\Re z<a$.
The target $\beta_m=m!u_m$ is the coefficient of the prescribed top pole.
If $u_m=0$, the actual pole order is lower.

\begin{theorem}[Direct recovery of angular logarithmic data]
\label{thm:vectorlogs}
Assume \eqref{eq:vectorclass} and observe
$Y(\nu)=G(a+i\nu)+e(\nu)$ with
$\sup_\nu\|e(\nu)\|_X\le\eta$.
For $\sigma,h,d>0$, define
\begin{align}
 E_{h,\sigma}(Y)&=\frac1{2\pi}\int_{\mathbb R}
 Y(\nu)h^{-a-i\nu}e^{\sigma^2(a+i\nu)^2/2}\,d\nu,
 \label{eq:vectoraverage}\\
 \widehat\beta_{m,h}(Y)&=d^{-m}\sum_{j=0}^m
 (-1)^{m-j}\binom mj E_{he^{-jd},\sigma}(Y).
 \label{eq:logestimator}
\end{align}
Set
\begin{equation}\label{eq:UVconstants}
 U=\frac{e^{\alpha^2\sigma^2/2}(1+e^{-\alpha d})^m}{d^m},
 \qquad
 V=\frac{e^{a^2\sigma^2/2}(1+e^{ad})^m}
 {\sqrt{2\pi}\,\sigma d^m}.
\end{equation}
Then
\begin{equation}\label{eq:vectorbound}
 \|\widehat\beta_{m,h}-\beta_m\|_X
 \le UMh^\alpha+V\eta h^{-a}.
\end{equation}
Every bounded linear functional $\ell\in X^*$ consequently satisfies
\begin{equation}\label{eq:chargeerror}
 |\ell(\widehat\beta_{m,h})-\ell(\beta_m)|
 \le\|\ell\|_{X^*}(UMh^\alpha+V\eta h^{-a}).
\end{equation}
For fixed $m,d,\sigma,a,\alpha$, optimization in $h$ gives an upper
bound proportional to
$M^{a/(a+\alpha)}\eta^{\alpha/(a+\alpha)}$.
No bound uniform in log degree $m$ is asserted.
\end{theorem}
\begin{proof}
The integrals converge absolutely in $X$, since
$\|G(a+i\nu)\|_X\le Q_a(f)$ and the frequency kernel is Gaussian.
Bochner Fubini and the scalar Gaussian integral give
\[
 E_{h,\sigma}(G)=\mathbb E f(he^{\sigma Z}),\qquad Z\sim N(0,1).
\]
Put $L=-\log h$ and
$S(t)=\mathbb E\sum_{k=0}^m u_k(t-\sigma Z)^k$.
This is an $X$-valued polynomial of degree at most $m$, with top
coefficient $u_m$. Its forward difference obeys
\[
 d^{-m}\sum_{j=0}^m(-1)^{m-j}\binom mj S(L+jd)=m!u_m.
\]
At scale $he^{-jd}$ the remainder contribution has norm at most
$Mh^\alpha e^{-\alpha jd}e^{\alpha^2\sigma^2/2}$.
The frequency-kernel $L^1$ norm is
$h^{-a}e^{ajd}e^{a^2\sigma^2/2}/(\sqrt{2\pi}\sigma)$.
Summing absolute values of the binomial weights proves
\eqref{eq:vectorbound}. Duality proves \eqref{eq:chargeerror}.
For $M,\eta>0$, the choice
\[
 h_* =\left(\frac{aV\eta}{\alpha UM}\right)^{1/(a+\alpha)}
\]
gives upper constant
\[
 U^{a/(a+\alpha)}V^{\alpha/(a+\alpha)}
 \left[(a/\alpha)^{\alpha/(a+\alpha)}
       +(\alpha/a)^{a/(a+\alpha)}\right].
\]
\end{proof}

\begin{theorem}[The fixed-degree error exponent cannot be improved]
\label{thm:loglower}
For $m\ge1$, let $\mathcal R_m(\eta)$ be the minimax norm error
for $\beta_m$ over the class \eqref{eq:vectorclass} and the stated
noise balls. Put $C_m=(m/(\alpha e))^m$ and
$\theta=\alpha/(a+\alpha)$. Then
\begin{equation}\label{eq:loglower}
 \mathcal R_m(\eta)\ge\frac12
 \left(\frac{m!M}{C_m}\right)^{a/(a+\alpha)}
 (2\eta a^{m+1})^\theta.
\end{equation}
Together with Theorem~\ref{thm:vectorlogs}, this proves the matching
minimax exponent for every fixed positive log degree. The constants
are not claimed to be sharp.
\end{theorem}
\begin{proof}
The zero-noise lower bound is zero and needs no construction; assume
$\eta>0$ in what follows.
Choose a unit vector $v\in X$ and consider
\[
 f_{h,c}(x)=c\,[\log(h/x)]_+^m v,
 \qquad P(t)=c(t+\log h)^m v,
\]
where $[y]_+=\max(y,0)$. For $x<h$ the remainder vanishes. For
$x\ge h$, set $s=\log(x/h)\ge0$. Then
\[
 \frac{\|f_{h,c}(x)-P(\log(1/x))\|_X}{x^\alpha}
 =ch^{-\alpha}s^me^{-\alpha s}\le ch^{-\alpha}C_m.
\]
The maximum occurs at $s=m/\alpha$.
Choose $c=Mh^\alpha/C_m$ to satisfy the same global prior.
Direct integration yields
\[
 G_{h,c}(z)=\frac{cm!h^z}{z^{m+1}}v,\qquad
 \|G_{h,c}\|_{a,\infty}=\frac{cm!h^a}{a^{m+1}}.
\]
Its top coefficient is $cm!v$. Set
\[
 h=\left(\frac{2\eta C_ma^{m+1}}{m!M}\right)^{1/(a+\alpha)}.
\]
Zero and $f_{h,c}$ have a common midpoint observation within
$\eta$ of each transform, and their target coefficients differ by
$cm!v$. The triangle inequality forces error at least $cm!/2$,
which is \eqref{eq:loglower}.
\end{proof}

This construction allows arbitrary lower polynomial coefficients,
as does the observation class. Additional known constraints on those
coefficients could change the minimax problem. For a nonzero functional
$\ell$, choose $v$ approaching its dual-norm supremum; the same argument
gives the charge-specific lower bound multiplied by $\|\ell\|$.
In a Hilbert space this factor is attained by a vector aligned with
the test function.

\begin{corollary}[Finite angular pairings and frequency bands]
\label{cor:angularband}
If $X=L^2(K,d\Omega)$ and $g\in L^2(K,d\Omega)$, define
\[
 Q_g=\int_K\overline{g(\zeta)}\,\beta_m(\zeta)\,d\Omega.
\]
Then \eqref{eq:chargeerror} holds with $\|\ell\|=\|g\|_2$.
For $X=H^{-s}(S^2)$ the corresponding bound uses the $H^s$ norm
of the dual test function. If data are available only for
$|\nu|\le\Omega$ and $Q_a(f)\le Q$ is known, truncate each
integral in \eqref{eq:logestimator}. Its norm error is at most
\begin{equation}\label{eq:logband}
 UMh^\alpha+Vh^{-a}\left[
 \eta\operatorname{erf}(\sigma\Omega/\sqrt2)
 +Q\operatorname{erfc}(\sigma\Omega/\sqrt2)\right].
\end{equation}
\end{corollary}
\begin{proof}
The first statements are Hilbert-space and Sobolev duality.
For each Gaussian integral the measured and omitted kernel norms
are its full $L^1$ norm times, respectively, the displayed error
function and complementary error function. Use $\|e\|\le\eta$
on the band and $\|G\|\le Q$ outside it, then sum binomial weights.
\end{proof}

For $m=0$, tuning the smoothing variance improves the explicit upper
constant. In the optimized bound its $\sigma$-dependent factor is
\[
 \exp\!\left(\frac{a\alpha\sigma^2}{2}\right)
 (\sqrt{2\pi}\sigma)^{-\alpha/(a+\alpha)}.
\]
Differentiating its logarithm gives the unique minimum
$\sigma^2=1/[a(a+\alpha)]$. Thus at $a=\alpha=1$,
\begin{equation}\label{eq:improvedconstant}
 \|\widehat\beta_0-\beta_0\|_X
 \le2(e/\pi)^{1/4}\sqrt{M\eta},\qquad
 h=\pi^{-1/4}\sqrt{\eta/M},\quad\sigma=1/\sqrt2.
\end{equation}
For the highest coefficient with $m>0$, the same variance minimizes this derived upper constant
at fixed $d$, because the remaining binomial factors do not depend
on $\sigma$. This optimizes the displayed bound, not all estimators.

\subsection{The entire principal part and the cost of lower coefficients}
\label{sec:fulljet}
The highest coefficient is only one part of a logarithmic soft
expansion. In a fixed reference-energy convention, write
$\beta_k=k!u_k$, $0\le k\le m$, for all coefficients of its Mellin
principal part. Lower coefficients can depend on the higher ones
under a change of convention, as Section~\ref{sec:units} explains.
We now give simultaneous estimators and componentwise optimal rates
for this full principal part. Holomorphic Laurent coefficients are
not recovery targets here.

At the same scales as Theorem~\ref{thm:vectorlogs}, put
$L=-\log h$ and define the Lagrange polynomials
\[
 \ell_j(t)=\prod_{\substack{0\le r\le m\\r\ne j}}
 \frac{t-L-rd}{d(j-r)},\qquad 0\le j\le m.
\]
For coefficient extraction denoted by $[t^n]$, set
\begin{equation}\label{eq:fullweights}
 W_{k,j}=\sum_{r=0}^{\lfloor(m-k)/2\rfloor}
 \frac{(-\sigma^2/2)^r}{r!}(k+2r)!
 [t^{k+2r}]\ell_j(t),\qquad
 \widehat\beta_k=\sum_{j=0}^m W_{k,j}E_{he^{-jd},\sigma}(Y).
\end{equation}
These are explicit finite real weights. No unknown soft coefficient
enters their construction. For $n\ge0$ define the positive polynomial
\[
 H_n(A;\sigma)=\sum_{r=0}^{\lfloor n/2\rfloor}
 \frac{(\sigma^2/2)^r A^{n-2r}}{r!(n-2r)!}.
\]

\begin{theorem}[Full principal-part recovery]\label{thm:fulljet}
Under \eqref{eq:vectorclass}, let $m\ge1$ and fix $d,\sigma>0$.
For every $k\le m$, the estimator \eqref{eq:fullweights} satisfies
\begin{align}
 \|\widehat\beta_k-\beta_k\|_X
 &\le \frac{H_{m-k}(|L|+md;\sigma)}{d^m}\,
 \mathcal E(h),\label{eq:fullbound}\\
 \mathcal E(h)&=
 e^{\alpha^2\sigma^2/2}Mh^\alpha(1+e^{-\alpha d})^m
 +\frac{e^{a^2\sigma^2/2}}{\sqrt{2\pi}\sigma}
 \eta h^{-a}(1+e^{ad})^m.\nonumber
\end{align}
Let $\mathcal R_{m,k}(\eta)$ be the minimax norm error for the
individual target $\beta_k$. As $\eta/M\downarrow0$, with all other
parameters fixed,
\begin{equation}\label{eq:fullrate}
 \mathcal R_{m,k}(\eta)\asymp
 M^{a/(a+\alpha)}\eta^{\alpha/(a+\alpha)}
 [\log(M/\eta)]^{m-k}.
\end{equation}
Here $\asymp$ means upper and lower positive constants independent of
$M,\eta$ in the stated small-noise regime. It does not assert a sharp
minimax constant or a bound uniform in $m$.
\end{theorem}
\begin{proof}
The polynomial component of each exact-data Gaussian average is
$S(t)=\mathbb E P(t-\sigma Z)$, where $P(t)=\sum u_kt^k$.
On degree-at-most-$m$ polynomials, Gaussian averaging is the finite
heat operator $S=e^{(\sigma^2/2)\partial_t^2}P$.
Lagrange interpolation reconstructs $S$ from $S(L+jd)$, and the
finite inverse operator $e^{-(\sigma^2/2)\partial_t^2}$ recovers $P$.
Applying $k![t^k]$ gives exactly \eqref{eq:fullweights}.
This argument uses the polynomial component of a convergent Gaussian
average; it does not assume that a global polynomial alone has an
absolutely convergent Mellin transform.

Put $A=|L|+md$. The denominator of $\ell_j$ has absolute value
$d^m j!(m-j)!$, and every root of its numerator has absolute value
at most $A$. Hence, for $0\le n\le m$,
\[
 |[t^n]\ell_j(t)|\le
 \frac{\binom mn A^{m-n}}{d^m j!(m-j)!}.
\]
Substitution in \eqref{eq:fullweights} gives
\[
 |W_{k,j}|\le\binom mj d^{-m}H_{m-k}(A;\sigma).
\]
Each observed Gaussian average differs from its polynomial component
by at most
\[
 e^{\alpha^2\sigma^2/2}Mh^\alpha e^{-\alpha jd}
 +\frac{e^{a^2\sigma^2/2}}{\sqrt{2\pi}\sigma}
 \eta h^{-a}e^{ajd}.
\]
Summation proves \eqref{eq:fullbound}.

For the upper rate choose
$h=(\eta/M)^{1/(a+\alpha)}$ and
$L=\log(M/\eta)/(a+\alpha)$.
Both terms in $\mathcal E(h)$ are a fixed constant times
$S_\eta=M^{a/(a+\alpha)}\eta^{\alpha/(a+\alpha)}$.
Since $H_{m-k}(|L|+md;\sigma)=O((1+|L|)^{m-k})$, the upper
bound in \eqref{eq:fullrate} follows.

For the lower rate use the same cutoff pair as
Theorem~\ref{thm:loglower}. Set
\[
 C_m=(m/(\alpha e))^m,\quad
 K_0=\frac{2a^{m+1}C_m}{m!},\quad
 h_0=(K_0\eta/M)^{1/(a+\alpha)},\quad L_0=-\log h_0.
\]
The coefficient of $t^k$ in
$P(t)=c(t-L_0)^m v$, with $c=Mh_0^\alpha/C_m$, is
$c\binom mk(-L_0)^{m-k}v$.
The common midpoint observation has noise at most $\eta$, so
\begin{equation}\label{eq:fulllower}
 \mathcal R_{m,k}(\eta)\ge
 \frac{1}{2(m-k)!}
 \left(\frac{m!}{C_m}\right)^{a/(a+\alpha)}
 (2a^{m+1})^{\alpha/(a+\alpha)}S_\eta |L_0|^{m-k}.
\end{equation}
Since $L_0=[\log(M/\eta)-\log K_0]/(a+\alpha)$, this proves
the matching small-noise lower rate. The zero-noise lower bound
is trivially zero.
\end{proof}

For example, at degree $m=2$, setting
$S_\eta=M^{a/(a+\alpha)}\eta^{\alpha/(a+\alpha)}$ gives
\begin{center}
\begin{tabular}{ll}
\toprule
Target coefficient & Componentwise minimax order as $\eta/M\downarrow0$\\
\midrule
Triple pole $\beta_2$ & $S_\eta$\\
Double pole $\beta_1$ & $S_\eta\log(M/\eta)$\\
Ordinary residue $\beta_0$ & $S_\eta[\log(M/\eta)]^2$\\
\bottomrule
\end{tabular}
\end{center}
This loss concerns unknown lower and higher coefficients in the stated
class, at fixed convention. Knowing some nuisance coefficients can
improve a different recovery problem. A bounded angular pairing
inherits each upper estimate with its dual norm and the corresponding
lower rate when the functional is nonzero. For a continuum frequency
band, the same weight bound applies after replacing the noise term
$\eta$ in $\mathcal E(h)$ by
$\eta\operatorname{erf}(\sigma\Omega/\sqrt2)
+Q\operatorname{erfc}(\sigma\Omega/\sqrt2)$.

Moving the interpolation window does not remove the lower bound for a
coefficient at a prescribed reference scale. Choosing a new reference
can instead change that lower coefficient, hence change the recovery
target. This is separate from a passive change of energy unit with a
co-transforming reference, which Section~\ref{sec:units} treats.


\subsection{Relation to physical soft charges and loop calculations}
The object $Q_g$ is a smeared coefficient. If an independent field-theory
derivation identifies a normalized physical charge as
\[
 Q_{\mathrm{phys}}[\varepsilon]
 =\ell_\varepsilon(\beta_m)+Q_{\mathrm{hard}}[\varepsilon],
\]
then the coefficient estimate transfers with factor
$\|\ell_\varepsilon\|$, plus any independently bounded hard-sector
error. For a differential angular map $D$, the dual test is
$D^*\varepsilon$, which must belong to $X^*$.
This transfer requires the charge map, normalization, boundary
conditions, and hard contribution from the physical theory.
The recovery theorem supplies none of those Ward-identity inputs.
The derivation of actual logarithmic graviton charges involves
radiative phase space and dressed fields \cite{wardlogs}.

Logarithmic soft terms and higher-pole celestial operators are established
physics \cite{sahoosen,loopcurrents,longrange}. The integer $m$ in
\eqref{eq:vectorclass} denotes log degree, not loop order: for example,
one-loop Yang--Mills soft currents can exhibit double and triple poles
\cite{loopcurrents}. Our statement quantifies recovery if a specified
fixed-order coefficient satisfies the stated uniform remainder bound.
It does not derive the coefficient from an arbitrary loop amplitude.

\subsection{Why a fixed-order estimate is not a resummation theorem}
\label{sec:resummation}
Two elementary Laplace--Mellin examples make the scope testable.
They illustrate the shifted poles and branch structures already
discussed in the celestial literature \cite{gu}.
Use an infrared kernel supported on $0<x\le1$ and set $t=\log(1/x)$.
For $\lambda>0$,
\begin{equation}\label{eq:shiftedpole}
 \int_0^\infty e^{-zt}e^{-\lambda t}\,dt=\frac1{z+\lambda},
 \qquad
 \int_0^\infty e^{-zt}\sum_{j=0}^n
 \frac{(-\lambda t)^j}{j!}\,dt
 =\sum_{j=0}^n\frac{(-\lambda)^j}{z^{j+1}}.
\end{equation}
Every finite truncation has poles at zero, while the resummed kernel
has a shifted pole. Their difference is not $O(e^{-\alpha t})$
as $t\to\infty$, so it fails the remainder prior of our theorem.

For $0<\gamma<1$, another exactly integrable model gives
\begin{align}
 H(z)&=\int_0^\infty e^{-zt}(1+\lambda t)^{-\gamma}\,dt
 \label{eq:branchmodel}\\
 &=\lambda^{-\gamma}z^{\gamma-1}e^{z/\lambda}
 \Gamma(1-\gamma,z/\lambda),\qquad \Re z>0,\nonumber\\
 z^{1-\gamma}H(z)&\longrightarrow
 \lambda^{-\gamma}\Gamma(1-\gamma)\quad(z\downarrow0).
 \label{eq:branchlimit}
\end{align}
The substitution $u=1+\lambda t$ proves the incomplete-gamma formula.
Alternatively, after $u=zt$, the last expression is
$\int_0^\infty e^{-u}(z+\lambda u)^{-\gamma}\,du$;
domination by $e^{-u}(\lambda u)^{-\gamma}$ proves the limit.
The fractional power corresponds to a branch point of the continuation,
not a finite Laurent pole. In contrast, truncating the formal binomial
series at degree $n$ gives
\[
 \sum_{j=0}^n\frac{(-1)^j(\gamma)_j\lambda^j}{z^{j+1}}.
\]
No interchange of an infinite perturbative sum and the Mellin integral
is used or justified here. These controls prevent a finite-log fit
from being reported as a proof of an all-orders resummed soft algebra.
